\paragraph{Current conditional multi-model adapter.}

The controlled experiment varies only the gold-price engine. Black--Scholes/
GBM, Heston, Bates--Poisson and Full Bates--Hawkes use the Team~8 parameters
recalibrated to the 2 September 2026 GLD surface. All engines start from the same
USD~4,417/oz level anchor, defined as Barrick's Q2 2026 average realised gold
price. It is neither a point-in-time spot observation nor a Team~4 simulated
price. Fourteen U.S. Treasury observations from 2 September 2026 are fitted
with the same-date Nelson--Siegel--Svensson curve (RMSE 2.0608 basis points
against the continuous par-yield proxy). Option calibration uses the fitted
rate at each maturity; quarterly paths use integral-preserving forward rates.
This is not a bootstrapped zero curve. The conditional bridge transfers option-implied shape under
$\mathbb Q$; it does not establish a one-for-one GLD conversion or a validated
$\mathbb Q$-to-$\mathbb P$ transformation.

This current NSS is independent of the legacy Team~4 demonstrator. The
operating contract is separate. Barrick's reported production and Cost of
Sales replace the first two quarters of 2026; the Team~4 production and cost
forecasts begin at 2026Q3. The legacy Team~4 NSS fit, implied-volatility inversion,
GBM paths and illustrative stochastic EBITDA valuation are explicitly
excluded. For model $j$, quarter $t$ and path $m$, the adapter forms
\begin{equation}
M_{t,j}^{(m)}=\frac{\bigl(P_{t,j}^{Au,(m)}-C_t\bigr)Q_t}{1000},
\end{equation}
where $Q_t$ is in thousand troy ounces and $C_t$ is in USD/oz. One common DCF,
WACC stream and equity bridge then maps every price branch to per-share value.

The run uses 8,192 paths, 260 fine steps aggregated to 20 quarters, aligned
price-engine seed 20260902 and independent WACC seed 20260902. The WACC shock
array is generated once and reused byte-identically. Therefore cross-model
differences arise from the Team~8 price law, while the actual/forecast
operating boundary and all corporate assumptions remain common.
